%
% claude-code-biff.tex
%
% Z Specification: Claude Code + Biff Workflow State (Layer 2)
% Extends the base Claude Code state machine with biff communication
% and workflow coordination state.
%
% Type-check: fuzz -t claude-code-biff.tex
% Animate:    probcli claude-code-biff.tex -animate 20
%

\documentclass[a4paper,10pt,fleqn]{article}
\usepackage[margin=1in]{geometry}
\usepackage{fuzz}
\usepackage[colorlinks=true,linkcolor=blue,citecolor=blue,urlcolor=blue]{hyperref}

\begin{document}

\title{Claude Code + Biff: A Z Specification}
\author{Layer 2 --- Workflow Coordination State}
\date{March 2026}
\hypersetup{
  pdftitle={Claude Code + Biff: A Z Specification},
  pdfauthor={Punt Labs},
  pdfsubject={Z Specification},
  pdfcreator={fuzz/probcli}
}
\maketitle

\tableofcontents
\newpage

% ============================================================
\section{Introduction}
% ============================================================

This specification extends the base Claude Code state machine
(\texttt{claude-code.tex}) with biff communication and workflow
coordination state. Biff provides team communication primitives
modeled after BSD Unix tools: \texttt{/plan}, \texttt{/who},
\texttt{/finger}, \texttt{/write}, \texttt{/read}, \texttt{/wall},
and \texttt{/talk}.

The key workflow invariants enforced by this model:

\begin{enumerate}
  \item \textbf{Plan before work}: Claude must set a plan (via
    \texttt{/plan} or auto-detect from git branch) before editing
    files or running commands.
  \item \textbf{Message awareness}: unread messages are surfaced
    before Claude declares ``done'' (Stop hook).
  \item \textbf{Wall awareness}: active wall broadcasts are injected
    into session context at startup.
  \item \textbf{Bead claim before work}: Claude must claim a bead
    (\texttt{bd update --status=in\_progress}) before editing files.
\end{enumerate}

% ============================================================
\section{Base Types (from claude-code.tex)}
% ============================================================

\begin{zed}
  [TOOLNAME, TOOLINPUT, TOOLOUTPUT, PROMPT, RESPONSE, AGENTID, SESSIONID, CONTEXTCHUNK]
\end{zed}

\begin{zed}
  SessionPhase ::= spInactive | spIdle | spProcessing | spResponding | spEnded
\end{zed}

\begin{zed}
  ToolPhase ::= tpNone | tpPreHook | tpPermission | tpExecuting | tpPostHook
\end{zed}

\begin{zed}
  PermissionMode ::= pmDefault | pmPlan | pmAcceptEdits | pmDontAsk | pmBypass
\end{zed}

\begin{zed}
  PermissionDecision ::= pdAllow | pdDeny | pdAsk
\end{zed}

\begin{zed}
  HookDecision ::= hdAllow | hdBlock
\end{zed}

\begin{zed}
  SessionStartSource ::= ssStartup | ssResume | ssClear | ssCompact
\end{zed}

\begin{zed}
  SessionEndReason ::= seClear | seLogout | sePromptExit | seBypassDisabled | seOther
\end{zed}

\begin{zed}
  ZBOOL ::= ztrue | zfalse
\end{zed}

% ============================================================
\section{Biff Types}
% ============================================================

\begin{zed}
  [BEADID, USERID, MSGID, WALLTEXT]
\end{zed}

\texttt{BEADID} identifies a tracked work item. \texttt{USERID}
identifies a team member. \texttt{MSGID} uniquely identifies a
message. \texttt{WALLTEXT} is the content of a wall broadcast.

\begin{zed}
  PlanSource ::= psManual | psAuto
\end{zed}

How the plan was set. \texttt{psManual} is an explicit \texttt{/plan}
command; \texttt{psAuto} is derived from the git branch name by the
\texttt{PostToolUse} Bash hook. A manual plan blocks auto-overwrite.

\begin{zed}
  ToolCategory ::= tcEdit | tcBash | tcOther
\end{zed}

Coarse tool classification for workflow gates. \texttt{tcEdit} covers
\texttt{Edit}, \texttt{Write}, and similar file-modifying tools.
\texttt{tcBash} covers shell commands. \texttt{tcOther} is everything
else (read-only tools, MCP queries, etc.).

% ============================================================
\section{Global Constants}
% ============================================================

\begin{axdef}
  maxTools : \nat \\
  maxSubagents : \nat \\
  maxContext : \nat \\
  maxUnread : \nat
  \where
  maxTools = 5 \\
  maxSubagents = 2 \\
  maxContext = 3 \\
  maxUnread = 3
\end{axdef}

% ============================================================
\section{State}
% ============================================================

The state combines the base Claude Code state with biff workflow
state in a single flat schema. Base invariants are preserved and
biff invariants are added.

\begin{schema}{State}
  %
  % --- Base Claude Code state ---
  %
  sessionId : \finset SESSIONID \\
  sessionPhase : SessionPhase \\
  permMode : PermissionMode \\
  toolPhase : ToolPhase \\
  currentTool : \finset TOOLNAME \\
  toolCategory : ToolCategory \\
  toolsThisTurn : \nat \\
  activeSubagents : \finset AGENTID \\
  context : \seq CONTEXTCHUNK \\
  pendingPrompt : \finset PROMPT \\
  pendingResponse : \finset RESPONSE \\
  %
  % --- Biff workflow state ---
  %
  planSet : ZBOOL \\
  planSource : \finset PlanSource \\
  beadClaimed : \finset BEADID \\
  unreadCount : \nat \\
  talkActive : ZBOOL \\
  wallActive : ZBOOL \\
  wallText : \finset WALLTEXT
  \where
  %
  % === Base invariants (from claude-code.tex) ===
  %
  \# sessionId \leq 1 \\
  sessionPhase = spInactive \implies sessionId = \emptyset \\
  sessionPhase \neq spInactive \implies \# sessionId = 1 \\
  sessionPhase \neq spProcessing \implies toolPhase = tpNone \\
  \# currentTool \leq 1 \\
  toolPhase = tpNone \implies currentTool = \emptyset \\
  toolPhase \neq tpNone \implies \# currentTool = 1 \\
  toolsThisTurn \leq maxTools \\
  sessionPhase = spInactive \implies activeSubagents = \emptyset \\
  sessionPhase = spEnded \implies activeSubagents = \emptyset \\
  \# activeSubagents \leq maxSubagents \\
  \# context \leq maxContext \\
  \# pendingPrompt \leq 1 \\
  sessionPhase = spIdle \implies pendingPrompt = \emptyset \\
  sessionPhase = spInactive \implies pendingPrompt = \emptyset \\
  sessionPhase = spEnded \implies pendingPrompt = \emptyset \\
  \# pendingResponse \leq 1 \\
  sessionPhase \neq spResponding \implies pendingResponse = \emptyset \\
  %
  % === Tool category tracks current tool ===
  %
  toolPhase = tpNone \implies toolCategory = tcOther \\
  %
  % === Biff invariants ===
  %
  % Plan source exists iff plan is set
  %
  \# planSource \leq 1 \\
  planSet = zfalse \implies planSource = \emptyset \\
  planSet = ztrue \implies \# planSource = 1 \\
  %
  % Bead claim: at most one active bead
  %
  \# beadClaimed \leq 1 \\
  %
  % Unread bounded for animation
  %
  unreadCount \leq maxUnread \\
  %
  % Wall text exists iff wall is active
  %
  \# wallText \leq 1 \\
  wallActive = zfalse \implies wallText = \emptyset \\
  wallActive = ztrue \implies \# wallText = 1 \\
  %
  % No biff state when session is inactive
  %
  sessionPhase = spInactive \implies planSet = zfalse \\
  sessionPhase = spInactive \implies beadClaimed = \emptyset \\
  sessionPhase = spInactive \implies unreadCount = 0 \\
  sessionPhase = spInactive \implies talkActive = zfalse \\
  sessionPhase = spInactive \implies wallActive = zfalse
\end{schema}

% ============================================================
\section{Initialization}
% ============================================================

\begin{schema}{Init}
  State'
  \where
  sessionId' = \emptyset \\
  sessionPhase' = spInactive \\
  permMode' = pmDefault \\
  toolPhase' = tpNone \\
  currentTool' = \emptyset \\
  toolCategory' = tcOther \\
  toolsThisTurn' = 0 \\
  activeSubagents' = \emptyset \\
  context' = \langle \rangle \\
  pendingPrompt' = \emptyset \\
  pendingResponse' = \emptyset \\
  %
  planSet' = zfalse \\
  planSource' = \emptyset \\
  beadClaimed' = \emptyset \\
  unreadCount' = 0 \\
  talkActive' = zfalse \\
  wallActive' = zfalse \\
  wallText' = \emptyset
\end{schema}

% ============================================================
\section{Biff Frame Condition}
% ============================================================

Most base operations do not change biff state. This abbreviation
schema captures ``biff state unchanged'' to reduce repetition.

\begin{schema}{BiffFrame}
  \Delta State
  \where
  planSet' = planSet \\
  planSource' = planSource \\
  beadClaimed' = beadClaimed \\
  unreadCount' = unreadCount \\
  talkActive' = talkActive \\
  wallActive' = wallActive \\
  wallText' = wallText
\end{schema}

% ============================================================
\section{Session Lifecycle Operations}
% ============================================================

\subsection{StartSession}

\texttt{SessionStart} hook fires. Biff's \texttt{session-start.sh}
initializes the session. If a wall is active, it is injected into
context.

\begin{schema}{StartSession}
  \Delta State \\
  sid? : SESSIONID \\
  source? : SessionStartSource \\
  mode? : PermissionMode \\
  hasWall? : ZBOOL \\
  wallContent? : \finset WALLTEXT
  \where
  sessionPhase = spInactive \\
  \# wallContent? \leq 1 \\
  hasWall? = zfalse \implies wallContent? = \emptyset \\
  hasWall? = ztrue \implies \# wallContent? = 1 \\
  %
  sessionPhase' = spIdle \\
  sessionId' = \{ sid? \} \\
  permMode' = mode? \\
  toolPhase' = tpNone \\
  currentTool' = \emptyset \\
  toolCategory' = tcOther \\
  toolsThisTurn' = 0 \\
  activeSubagents' = \emptyset \\
  context' = \langle \rangle \\
  pendingPrompt' = \emptyset \\
  pendingResponse' = \emptyset \\
  %
  planSet' = zfalse \\
  planSource' = \emptyset \\
  beadClaimed' = \emptyset \\
  unreadCount' = 0 \\
  talkActive' = zfalse \\
  wallActive' = hasWall? \\
  wallText' = wallContent?
\end{schema}

\subsection{SubmitPrompt}

\begin{schema}{SubmitPrompt}
  BiffFrame \\
  prompt? : PROMPT \\
  hookResult? : HookDecision \\
  contextChunk? : CONTEXTCHUNK
  \where
  sessionPhase = spIdle \\
  hookResult? = hdAllow \\
  sessionPhase' = spProcessing \\
  pendingPrompt' = \{ prompt? \} \\
  context' = (1 \upto maxContext) \dres (context \cat \langle contextChunk? \rangle) \\
  toolsThisTurn' = 0 \\
  toolPhase' = tpNone \\
  currentTool' = \emptyset \\
  toolCategory' = tcOther \\
  sessionId' = sessionId \\
  permMode' = permMode \\
  activeSubagents' = activeSubagents \\
  pendingResponse' = \emptyset
\end{schema}

\begin{schema}{RejectPrompt}
  \Xi State \\
  prompt? : PROMPT \\
  hookResult? : HookDecision
  \where
  sessionPhase = spIdle \\
  hookResult? = hdBlock
\end{schema}

% ============================================================
\section{Tool Call Operations}
% ============================================================

\subsection{BeginToolCall}

Claude decides to invoke a tool. The tool category is classified.

\begin{schema}{BeginToolCall}
  BiffFrame \\
  tool? : TOOLNAME \\
  category? : ToolCategory
  \where
  sessionPhase = spProcessing \\
  toolPhase = tpNone \\
  toolsThisTurn < maxTools \\
  %
  toolPhase' = tpPreHook \\
  currentTool' = \{ tool? \} \\
  toolCategory' = category? \\
  %
  sessionPhase' = sessionPhase \\
  sessionId' = sessionId \\
  permMode' = permMode \\
  toolsThisTurn' = toolsThisTurn \\
  activeSubagents' = activeSubagents \\
  context' = context \\
  pendingPrompt' = pendingPrompt \\
  pendingResponse' = pendingResponse
\end{schema}

\subsection{PreToolHookAllow --- with workflow gate}

The \texttt{PreToolUse} hook checks workflow prerequisites for
edit and bash tools. The hook allows only if the plan is set and
a bead is claimed (for edit/bash tools), or unconditionally for
other tools.

\begin{schema}{PreToolHookAllow}
  BiffFrame \\
  hookDecision? : PermissionDecision \\
  needsPermission? : ZBOOL
  \where
  sessionPhase = spProcessing \\
  toolPhase = tpPreHook \\
  hookDecision? = pdAllow \\
  %
  % WORKFLOW GATE: edit/bash tools require plan + bead
  %
  toolCategory \in \{ tcEdit, tcBash \} \implies planSet = ztrue \\
  toolCategory \in \{ tcEdit, tcBash \} \implies beadClaimed \neq \emptyset \\
  %
  toolPhase' = \IF needsPermission? = ztrue \THEN tpPermission \ELSE tpExecuting \\
  %
  sessionPhase' = sessionPhase \\
  sessionId' = sessionId \\
  permMode' = permMode \\
  currentTool' = currentTool \\
  toolCategory' = toolCategory \\
  toolsThisTurn' = toolsThisTurn \\
  activeSubagents' = activeSubagents \\
  context' = context \\
  pendingPrompt' = pendingPrompt \\
  pendingResponse' = pendingResponse
\end{schema}

\subsection{PreToolHookDeny}

The hook denies the tool call (workflow prerequisites not met,
or explicit deny for other reasons).

\begin{schema}{PreToolHookDeny}
  BiffFrame \\
  hookDecision? : PermissionDecision \\
  denialContext? : CONTEXTCHUNK
  \where
  sessionPhase = spProcessing \\
  toolPhase = tpPreHook \\
  hookDecision? = pdDeny \\
  %
  toolPhase' = tpNone \\
  currentTool' = \emptyset \\
  toolCategory' = tcOther \\
  context' = (1 \upto maxContext) \dres (context \cat \langle denialContext? \rangle) \\
  %
  sessionPhase' = sessionPhase \\
  sessionId' = sessionId \\
  permMode' = permMode \\
  toolsThisTurn' = toolsThisTurn \\
  activeSubagents' = activeSubagents \\
  pendingPrompt' = pendingPrompt \\
  pendingResponse' = pendingResponse
\end{schema}

\subsection{PreToolHookEscalate}

\begin{schema}{PreToolHookEscalate}
  BiffFrame \\
  hookDecision? : PermissionDecision
  \where
  sessionPhase = spProcessing \\
  toolPhase = tpPreHook \\
  hookDecision? = pdAsk \\
  %
  toolPhase' = tpPermission \\
  sessionPhase' = sessionPhase \\
  sessionId' = sessionId \\
  permMode' = permMode \\
  currentTool' = currentTool \\
  toolCategory' = toolCategory \\
  toolsThisTurn' = toolsThisTurn \\
  activeSubagents' = activeSubagents \\
  context' = context \\
  pendingPrompt' = pendingPrompt \\
  pendingResponse' = pendingResponse
\end{schema}

\subsection{GrantPermission}

\begin{schema}{GrantPermission}
  BiffFrame
  \where
  sessionPhase = spProcessing \\
  toolPhase = tpPermission \\
  toolPhase' = tpExecuting \\
  sessionPhase' = sessionPhase \\
  sessionId' = sessionId \\
  permMode' = permMode \\
  currentTool' = currentTool \\
  toolCategory' = toolCategory \\
  toolsThisTurn' = toolsThisTurn \\
  activeSubagents' = activeSubagents \\
  context' = context \\
  pendingPrompt' = pendingPrompt \\
  pendingResponse' = pendingResponse
\end{schema}

\subsection{DenyPermission}

\begin{schema}{DenyPermission}
  BiffFrame \\
  denialContext? : CONTEXTCHUNK
  \where
  sessionPhase = spProcessing \\
  toolPhase = tpPermission \\
  toolPhase' = tpNone \\
  currentTool' = \emptyset \\
  toolCategory' = tcOther \\
  context' = (1 \upto maxContext) \dres (context \cat \langle denialContext? \rangle) \\
  sessionPhase' = sessionPhase \\
  sessionId' = sessionId \\
  permMode' = permMode \\
  toolsThisTurn' = toolsThisTurn \\
  activeSubagents' = activeSubagents \\
  pendingPrompt' = pendingPrompt \\
  pendingResponse' = pendingResponse
\end{schema}

\subsection{CompleteToolSuccess}

\begin{schema}{CompleteToolSuccess}
  BiffFrame \\
  output? : CONTEXTCHUNK
  \where
  sessionPhase = spProcessing \\
  toolPhase = tpExecuting \\
  toolPhase' = tpNone \\
  currentTool' = \emptyset \\
  toolCategory' = tcOther \\
  toolsThisTurn' = toolsThisTurn + 1 \\
  context' = (1 \upto maxContext) \dres (context \cat \langle output? \rangle) \\
  sessionPhase' = sessionPhase \\
  sessionId' = sessionId \\
  permMode' = permMode \\
  activeSubagents' = activeSubagents \\
  pendingPrompt' = pendingPrompt \\
  pendingResponse' = pendingResponse
\end{schema}

\subsection{CompleteToolFailure}

\begin{schema}{CompleteToolFailure}
  BiffFrame \\
  errorContext? : CONTEXTCHUNK
  \where
  sessionPhase = spProcessing \\
  toolPhase = tpExecuting \\
  toolPhase' = tpNone \\
  currentTool' = \emptyset \\
  toolCategory' = tcOther \\
  toolsThisTurn' = toolsThisTurn + 1 \\
  context' = (1 \upto maxContext) \dres (context \cat \langle errorContext? \rangle) \\
  sessionPhase' = sessionPhase \\
  sessionId' = sessionId \\
  permMode' = permMode \\
  activeSubagents' = activeSubagents \\
  pendingPrompt' = pendingPrompt \\
  pendingResponse' = pendingResponse
\end{schema}

% ============================================================
\section{Biff Workflow Operations}
% ============================================================

These operations model biff tool calls that change workflow state.
They occur during \texttt{spProcessing} as tool completions.

\subsection{SetPlanManual}

The user or Claude explicitly calls \texttt{/plan}. This sets the
plan and blocks auto-overwrite.

\begin{schema}{SetPlanManual}
  \Delta State
  \where
  sessionPhase = spProcessing \\
  toolPhase = tpNone \\
  %
  planSet' = ztrue \\
  planSource' = \{ psManual \} \\
  %
  % Frame: everything else unchanged
  %
  sessionPhase' = sessionPhase \\
  sessionId' = sessionId \\
  permMode' = permMode \\
  toolPhase' = toolPhase \\
  currentTool' = currentTool \\
  toolCategory' = toolCategory \\
  toolsThisTurn' = toolsThisTurn \\
  activeSubagents' = activeSubagents \\
  context' = context \\
  pendingPrompt' = pendingPrompt \\
  pendingResponse' = pendingResponse \\
  beadClaimed' = beadClaimed \\
  unreadCount' = unreadCount \\
  talkActive' = talkActive \\
  wallActive' = wallActive \\
  wallText' = wallText
\end{schema}

\subsection{SetPlanAuto}

The \texttt{PostToolUse} Bash hook detects a git branch and
auto-sets the plan. Only succeeds if no manual plan is set.

\begin{schema}{SetPlanAuto}
  \Delta State
  \where
  sessionPhase = spProcessing \\
  toolPhase = tpNone \\
  planSource \neq \{ psManual \} \\
  %
  planSet' = ztrue \\
  planSource' = \{ psAuto \} \\
  %
  sessionPhase' = sessionPhase \\
  sessionId' = sessionId \\
  permMode' = permMode \\
  toolPhase' = toolPhase \\
  currentTool' = currentTool \\
  toolCategory' = toolCategory \\
  toolsThisTurn' = toolsThisTurn \\
  activeSubagents' = activeSubagents \\
  context' = context \\
  pendingPrompt' = pendingPrompt \\
  pendingResponse' = pendingResponse \\
  beadClaimed' = beadClaimed \\
  unreadCount' = unreadCount \\
  talkActive' = talkActive \\
  wallActive' = wallActive \\
  wallText' = wallText
\end{schema}

\subsection{ClaimBead}

Claude runs \texttt{bd update <id> --status=in\_progress}.

\begin{schema}{ClaimBead}
  \Delta State \\
  bead? : BEADID
  \where
  sessionPhase = spProcessing \\
  toolPhase = tpNone \\
  %
  beadClaimed' = \{ bead? \} \\
  %
  sessionPhase' = sessionPhase \\
  sessionId' = sessionId \\
  permMode' = permMode \\
  toolPhase' = toolPhase \\
  currentTool' = currentTool \\
  toolCategory' = toolCategory \\
  toolsThisTurn' = toolsThisTurn \\
  activeSubagents' = activeSubagents \\
  context' = context \\
  pendingPrompt' = pendingPrompt \\
  pendingResponse' = pendingResponse \\
  planSet' = planSet \\
  planSource' = planSource \\
  unreadCount' = unreadCount \\
  talkActive' = talkActive \\
  wallActive' = wallActive \\
  wallText' = wallText
\end{schema}

\subsection{CloseBead}

Claude runs \texttt{bd close <id>}.

\begin{schema}{CloseBead}
  \Delta State \\
  bead? : BEADID
  \where
  sessionPhase = spProcessing \\
  toolPhase = tpNone \\
  bead? \in beadClaimed \\
  %
  beadClaimed' = beadClaimed \setminus \{ bead? \} \\
  %
  sessionPhase' = sessionPhase \\
  sessionId' = sessionId \\
  permMode' = permMode \\
  toolPhase' = toolPhase \\
  currentTool' = currentTool \\
  toolCategory' = toolCategory \\
  toolsThisTurn' = toolsThisTurn \\
  activeSubagents' = activeSubagents \\
  context' = context \\
  pendingPrompt' = pendingPrompt \\
  pendingResponse' = pendingResponse \\
  planSet' = planSet \\
  planSource' = planSource \\
  unreadCount' = unreadCount \\
  talkActive' = talkActive \\
  wallActive' = wallActive \\
  wallText' = wallText
\end{schema}

\subsection{ReceiveMessage}

An external message arrives (NATS notification). Non-blocking,
increments unread count.

\begin{schema}{ReceiveMessage}
  \Delta State
  \where
  sessionPhase \neq spInactive \\
  sessionPhase \neq spEnded \\
  unreadCount < maxUnread \\
  %
  unreadCount' = unreadCount + 1 \\
  %
  sessionPhase' = sessionPhase \\
  sessionId' = sessionId \\
  permMode' = permMode \\
  toolPhase' = toolPhase \\
  currentTool' = currentTool \\
  toolCategory' = toolCategory \\
  toolsThisTurn' = toolsThisTurn \\
  activeSubagents' = activeSubagents \\
  context' = context \\
  pendingPrompt' = pendingPrompt \\
  pendingResponse' = pendingResponse \\
  planSet' = planSet \\
  planSource' = planSource \\
  beadClaimed' = beadClaimed \\
  talkActive' = talkActive \\
  wallActive' = wallActive \\
  wallText' = wallText
\end{schema}

\subsection{ReadMessages}

Claude calls \texttt{/read}. Clears unread count.

\begin{schema}{ReadMessages}
  \Delta State
  \where
  sessionPhase = spProcessing \\
  toolPhase = tpNone \\
  %
  unreadCount' = 0 \\
  %
  sessionPhase' = sessionPhase \\
  sessionId' = sessionId \\
  permMode' = permMode \\
  toolPhase' = toolPhase \\
  currentTool' = currentTool \\
  toolCategory' = toolCategory \\
  toolsThisTurn' = toolsThisTurn \\
  activeSubagents' = activeSubagents \\
  context' = context \\
  pendingPrompt' = pendingPrompt \\
  pendingResponse' = pendingResponse \\
  planSet' = planSet \\
  planSource' = planSource \\
  beadClaimed' = beadClaimed \\
  talkActive' = talkActive \\
  wallActive' = wallActive \\
  wallText' = wallText
\end{schema}

\subsection{StartTalk}

Claude calls \texttt{/talk @user}.

\begin{schema}{StartTalk}
  \Delta State
  \where
  sessionPhase = spProcessing \\
  toolPhase = tpNone \\
  talkActive = zfalse \\
  %
  talkActive' = ztrue \\
  %
  sessionPhase' = sessionPhase \\
  sessionId' = sessionId \\
  permMode' = permMode \\
  toolPhase' = toolPhase \\
  currentTool' = currentTool \\
  toolCategory' = toolCategory \\
  toolsThisTurn' = toolsThisTurn \\
  activeSubagents' = activeSubagents \\
  context' = context \\
  pendingPrompt' = pendingPrompt \\
  pendingResponse' = pendingResponse \\
  planSet' = planSet \\
  planSource' = planSource \\
  beadClaimed' = beadClaimed \\
  unreadCount' = unreadCount \\
  wallActive' = wallActive \\
  wallText' = wallText
\end{schema}

\subsection{EndTalk}

Claude calls \texttt{/talk\_end}.

\begin{schema}{EndTalk}
  \Delta State
  \where
  sessionPhase = spProcessing \\
  toolPhase = tpNone \\
  talkActive = ztrue \\
  %
  talkActive' = zfalse \\
  %
  sessionPhase' = sessionPhase \\
  sessionId' = sessionId \\
  permMode' = permMode \\
  toolPhase' = toolPhase \\
  currentTool' = currentTool \\
  toolCategory' = toolCategory \\
  toolsThisTurn' = toolsThisTurn \\
  activeSubagents' = activeSubagents \\
  context' = context \\
  pendingPrompt' = pendingPrompt \\
  pendingResponse' = pendingResponse \\
  planSet' = planSet \\
  planSource' = planSource \\
  beadClaimed' = beadClaimed \\
  unreadCount' = unreadCount \\
  wallActive' = wallActive \\
  wallText' = wallText
\end{schema}

\subsection{PostWall}

Claude calls \texttt{/wall <message>}.

\begin{schema}{PostWall}
  \Delta State \\
  text? : WALLTEXT
  \where
  sessionPhase = spProcessing \\
  toolPhase = tpNone \\
  %
  wallActive' = ztrue \\
  wallText' = \{ text? \} \\
  %
  sessionPhase' = sessionPhase \\
  sessionId' = sessionId \\
  permMode' = permMode \\
  toolPhase' = toolPhase \\
  currentTool' = currentTool \\
  toolCategory' = toolCategory \\
  toolsThisTurn' = toolsThisTurn \\
  activeSubagents' = activeSubagents \\
  context' = context \\
  pendingPrompt' = pendingPrompt \\
  pendingResponse' = pendingResponse \\
  planSet' = planSet \\
  planSource' = planSource \\
  beadClaimed' = beadClaimed \\
  unreadCount' = unreadCount \\
  talkActive' = talkActive
\end{schema}

\subsection{ClearWall}

Claude calls \texttt{/wall --clear}.

\begin{schema}{ClearWall}
  \Delta State
  \where
  sessionPhase = spProcessing \\
  toolPhase = tpNone \\
  wallActive = ztrue \\
  %
  wallActive' = zfalse \\
  wallText' = \emptyset \\
  %
  sessionPhase' = sessionPhase \\
  sessionId' = sessionId \\
  permMode' = permMode \\
  toolPhase' = toolPhase \\
  currentTool' = currentTool \\
  toolCategory' = toolCategory \\
  toolsThisTurn' = toolsThisTurn \\
  activeSubagents' = activeSubagents \\
  context' = context \\
  pendingPrompt' = pendingPrompt \\
  pendingResponse' = pendingResponse \\
  planSet' = planSet \\
  planSource' = planSource \\
  beadClaimed' = beadClaimed \\
  unreadCount' = unreadCount \\
  talkActive' = talkActive
\end{schema}

% ============================================================
\section{Subagent Operations}
% ============================================================

\subsection{SpawnSubagent}

\begin{schema}{SpawnSubagent}
  BiffFrame \\
  agentId? : AGENTID
  \where
  sessionPhase = spProcessing \\
  agentId? \notin activeSubagents \\
  \# activeSubagents < maxSubagents \\
  activeSubagents' = activeSubagents \cup \{ agentId? \} \\
  sessionPhase' = sessionPhase \\
  sessionId' = sessionId \\
  permMode' = permMode \\
  toolPhase' = toolPhase \\
  currentTool' = currentTool \\
  toolCategory' = toolCategory \\
  toolsThisTurn' = toolsThisTurn \\
  context' = context \\
  pendingPrompt' = pendingPrompt \\
  pendingResponse' = pendingResponse
\end{schema}

\subsection{StopSubagent}

\begin{schema}{StopSubagent}
  BiffFrame \\
  agentId? : AGENTID \\
  hookResult? : HookDecision \\
  resultContext? : CONTEXTCHUNK
  \where
  sessionPhase = spProcessing \\
  agentId? \in activeSubagents \\
  hookResult? = hdAllow \\
  activeSubagents' = activeSubagents \setminus \{ agentId? \} \\
  context' = (1 \upto maxContext) \dres (context \cat \langle resultContext? \rangle) \\
  sessionPhase' = sessionPhase \\
  sessionId' = sessionId \\
  permMode' = permMode \\
  toolPhase' = toolPhase \\
  currentTool' = currentTool \\
  toolCategory' = toolCategory \\
  toolsThisTurn' = toolsThisTurn \\
  pendingPrompt' = pendingPrompt \\
  pendingResponse' = pendingResponse
\end{schema}

% ============================================================
\section{Response and Stop Operations}
% ============================================================

\subsection{BeginResponse}

\begin{schema}{BeginResponse}
  BiffFrame \\
  response? : RESPONSE
  \where
  sessionPhase = spProcessing \\
  toolPhase = tpNone \\
  activeSubagents = \emptyset \\
  sessionPhase' = spResponding \\
  pendingResponse' = \{ response? \} \\
  pendingPrompt' = \emptyset \\
  sessionId' = sessionId \\
  permMode' = permMode \\
  toolPhase' = tpNone \\
  currentTool' = \emptyset \\
  toolCategory' = tcOther \\
  toolsThisTurn' = toolsThisTurn \\
  activeSubagents' = activeSubagents \\
  context' = context
\end{schema}

\subsection{FinishResponse --- with message awareness}

The \texttt{Stop} hook checks for unread messages. If unread
messages exist, the hook injects a reminder into context but
does \emph{not} block --- this is a soft gate (the agent should
read messages but is not forced to).

\begin{schema}{FinishResponse}
  BiffFrame \\
  hookResult? : HookDecision \\
  responseContext? : CONTEXTCHUNK
  \where
  sessionPhase = spResponding \\
  hookResult? = hdAllow \\
  sessionPhase' = spIdle \\
  pendingResponse' = \emptyset \\
  context' = (1 \upto maxContext) \dres (context \cat \langle responseContext? \rangle) \\
  toolsThisTurn' = 0 \\
  sessionId' = sessionId \\
  permMode' = permMode \\
  toolPhase' = tpNone \\
  currentTool' = \emptyset \\
  toolCategory' = tcOther \\
  activeSubagents' = activeSubagents \\
  pendingPrompt' = \emptyset
\end{schema}

\subsection{ForceContinue}

\begin{schema}{ForceContinue}
  BiffFrame \\
  hookResult? : HookDecision
  \where
  sessionPhase = spResponding \\
  hookResult? = hdBlock \\
  sessionPhase' = spProcessing \\
  pendingResponse' = \emptyset \\
  sessionId' = sessionId \\
  permMode' = permMode \\
  toolPhase' = tpNone \\
  currentTool' = \emptyset \\
  toolCategory' = tcOther \\
  toolsThisTurn' = toolsThisTurn \\
  activeSubagents' = activeSubagents \\
  context' = context \\
  pendingPrompt' = pendingPrompt
\end{schema}

% ============================================================
\section{Context Management}
% ============================================================

\subsection{CompactContext}

\begin{schema}{CompactContext}
  BiffFrame \\
  newLen? : \nat
  \where
  sessionPhase = spIdle \\
  \# context > 0 \\
  newLen? < \# context \\
  newLen? \leq maxContext \\
  context' = (1 \upto newLen?) \dres context \\
  sessionPhase' = sessionPhase \\
  sessionId' = sessionId \\
  permMode' = permMode \\
  toolPhase' = toolPhase \\
  currentTool' = currentTool \\
  toolCategory' = toolCategory \\
  toolsThisTurn' = toolsThisTurn \\
  activeSubagents' = activeSubagents \\
  pendingPrompt' = pendingPrompt \\
  pendingResponse' = pendingResponse
\end{schema}

% ============================================================
\section{Session End}
% ============================================================

\subsection{EndSession}

\begin{schema}{EndSession}
  \Delta State \\
  reason? : SessionEndReason
  \where
  sessionPhase = spIdle \\
  activeSubagents = \emptyset \\
  %
  sessionPhase' = spEnded \\
  sessionId' = \emptyset \\
  toolPhase' = tpNone \\
  currentTool' = \emptyset \\
  toolCategory' = tcOther \\
  toolsThisTurn' = 0 \\
  activeSubagents' = \emptyset \\
  context' = \langle \rangle \\
  pendingPrompt' = \emptyset \\
  pendingResponse' = \emptyset \\
  permMode' = permMode \\
  %
  % Biff state cleared
  %
  planSet' = zfalse \\
  planSource' = \emptyset \\
  beadClaimed' = \emptyset \\
  unreadCount' = 0 \\
  talkActive' = zfalse \\
  wallActive' = zfalse \\
  wallText' = \emptyset
\end{schema}

% ============================================================
\section{System Invariants}
% ============================================================

The specification enforces these properties across all reachable states:

\begin{enumerate}
  \item \textbf{Plan before editing} (workflow gate): the
    \texttt{PreToolHookAllow} operation requires \texttt{planSet = ztrue}
    and \texttt{beadClaimed $\neq \emptyset$} for edit and bash tools.
    No path through the state machine can reach tool execution for
    file-modifying tools without these prerequisites.

  \item \textbf{Manual plan blocks auto}: \texttt{SetPlanAuto} has
    precondition \texttt{planSource $\neq$ \{psManual\}}, preventing
    auto-detection from overwriting an explicit plan.

  \item \textbf{Wall consistency}: wall text exists if and only if
    the wall is active. No orphan text, no active wall without content.

  \item \textbf{Plan consistency}: plan source exists if and only if
    plan is set.

  \item \textbf{Clean session boundaries}: all biff state is cleared
    on \texttt{Init} and \texttt{EndSession}. No state leaks between
    sessions.

  \item \textbf{Talk exclusivity}: \texttt{StartTalk} requires
    \texttt{talkActive = zfalse}; \texttt{EndTalk} requires
    \texttt{talkActive = ztrue}. No double-start or double-end.

  \item \textbf{All base invariants preserved}: tool phase, session
    phase, subagent, and context invariants from the base model hold.
\end{enumerate}

% ============================================================
\section{Hook Event Map}
% ============================================================

\begin{tabular}{lllll}
\hline
\textbf{Hook} & \textbf{Matcher} & \textbf{Operation} & \textbf{Biff Effect} & \textbf{Blocks?} \\
\hline
SessionStart & * & StartSession & Load wall & No \\
PostToolUse & Bash & SetPlanAuto & Auto plan from branch & No \\
PreToolUse & Edit|Write & PreToolHookAllow/Deny & Require plan+bead & Yes \\
PreToolUse & Bash & PreToolHookAllow/Deny & Require plan+bead & Yes \\
PostToolUse & biff\_\_plan & SetPlanManual & Set manual plan & No \\
PostToolUse & biff\_\_write & (side effect) & Message sent & No \\
PostToolUse & biff\_\_read & ReadMessages & Clear unread & No \\
PostToolUse & biff\_\_talk & StartTalk & Enter talk mode & No \\
PostToolUse & biff\_\_talk\_end & EndTalk & Exit talk mode & No \\
PostToolUse & biff\_\_wall & PostWall/ClearWall & Update wall & No \\
Stop & * & FinishResponse & Check unread & Yes \\
SessionEnd & * & EndSession & Clear all & No \\
\hline
\end{tabular}

% ============================================================
\section{Use Case Integration Validation}
% ============================================================

This section cross-references each use case from
\texttt{biff\_coordination\_use\_cases\_v0.2.md} against the Z
specification operations and state variables. For each use case,
we identify: (1) which Z operations model the behavior, (2) what
the model checker proves about the integration path, and (3) any
gaps where the spec does not yet cover the use case.

% ------------------------------------------------------------
\subsection{Group A: Presence}
% ------------------------------------------------------------

\subsubsection{UC1 --- Check Who's Working on What}

\textbf{Z coverage:} Not modeled as an operation. \texttt{/who} is a
read-only query that does not change state. However, the state it reads
(\texttt{planSet}, \texttt{planSource}, \texttt{talkActive}) is maintained
by the model.

\textbf{What the model proves:} The \texttt{planSet} and
\texttt{planSource} variables are always consistent (invariant 4:
plan source exists iff plan is set). Any \texttt{/who} query will
see coherent plan state --- never a dangling source without a plan,
or a plan without a source.

\textbf{Gap:} None for V1. The idle time display relies on
\texttt{last\_active} timestamps in the relay, which is outside the
model's scope (relay internals are opaque).

\subsubsection{UC2 --- Inspect a Teammate's Availability}

\textbf{Z coverage:} Not modeled as an operation. Same rationale as UC1
--- read-only query. The \texttt{talkActive} state is accurate because
\texttt{StartTalk} and \texttt{EndTalk} enforce exclusivity (invariant 6).

\textbf{What the model proves:} A \texttt{/finger} query will never show
a session simultaneously in talk and not-in-talk. The talk state is
consistent across all reachable states.

\textbf{Gap:} None. The \texttt{mesgMode} (DND) variable was proposed
but not included in this model. If DND status needs formal verification,
add \texttt{mesgMode : ZBOOL} to the state schema.

\subsubsection{UC3 --- Declare What You're Working On}

\textbf{Z coverage:} Fully modeled.

\begin{itemize}
  \item Manual path: \texttt{SetPlanManual} --- sets \texttt{planSet = ztrue},
    \texttt{planSource = \{psManual\}}.
  \item Auto path: \texttt{SetPlanAuto} --- sets \texttt{planSet = ztrue},
    \texttt{planSource = \{psAuto\}}, guarded by
    \texttt{planSource $\neq$ \{psManual\}}.
  \item \texttt{/tty} naming is outside the model (relay metadata, no
    workflow impact).
\end{itemize}

\textbf{What the model proves:}

\begin{enumerate}
  \item \textbf{Manual blocks auto}: \texttt{SetPlanAuto} has precondition
    \texttt{planSource $\neq$ \{psManual\}}. The model checker verified
    across 161K states that no auto-plan overwrites a manual plan.
  \item \textbf{Plan is prerequisite for work}: Once the plan-before-work
    gate is active, the model proves that no file edit or bash command
    executes without \texttt{planSet = ztrue}. This validates the use
    case's claim that plan declaration is meaningful --- it is enforced,
    not advisory.
\end{enumerate}

\textbf{Gap:} Bead ID expansion (step 3 in the use case) is a relay-side
detail, not modeled.

% ------------------------------------------------------------
\subsection{Group B: Asynchronous Messaging}
% ------------------------------------------------------------

\subsubsection{UC4 --- Send a Focused Message}

\textbf{Z coverage:} Partially modeled. The send action is a tool
completion side effect (noted in the Hook Event Map as
\texttt{PostToolUse} on \texttt{biff\_\_write}). The model does not
track outbound messages because they do not affect the sender's
workflow state.

\textbf{What the model proves:} Sending a message does not disrupt
any workflow invariant. The \texttt{BiffFrame} ensures all biff state
is preserved through tool completions. No state corruption from
interleaving \texttt{/write} with other operations.

\textbf{Gap:} Recipient-side effects (unread count increment, push
notification) are modeled by \texttt{ReceiveMessage}, but the causal
link between send and receive is not captured --- they are independent
events in the model, which matches reality (async, fire-and-forget).

\subsubsection{UC5 --- Check and Process Incoming Messages}

\textbf{Z coverage:} Fully modeled.

\begin{itemize}
  \item \texttt{ReceiveMessage} --- increments \texttt{unreadCount} when
    a message arrives (external event, can fire in any active session phase).
  \item \texttt{ReadMessages} --- resets \texttt{unreadCount} to 0 when
    the agent calls \texttt{/read}.
\end{itemize}

\textbf{What the model proves:}

\begin{enumerate}
  \item \texttt{unreadCount} is bounded (\texttt{$\leq$ maxUnread}), so
    the model is finite and fully explorable.
  \item \texttt{ReceiveMessage} can fire during any session phase
    (idle, processing, responding) except \texttt{spInactive} and
    \texttt{spEnded}. This correctly models NATS push notifications
    arriving asynchronously.
  \item \texttt{ReadMessages} requires \texttt{spProcessing} and
    \texttt{toolPhase = tpNone}, which correctly reflects that \texttt{/read}
    is a tool call during the agentic loop.
\end{enumerate}

\textbf{Gap:} POP semantics (messages consumed once, then deleted) are
not modeled --- the count goes to 0, but we do not track individual
message identity. Sufficient for workflow invariants.

\subsubsection{UC7 --- Escalate a Decision to a Human}

\textbf{Z coverage:} Not modeled as a distinct operation. The use case
document proposes a \texttt{Stop} hook that injects escalation guidance.
In the Z spec, the \texttt{FinishResponse} operation's \texttt{Stop}
hook is the attachment point.

\textbf{What the model proves:} The \texttt{Stop} hook fires on every
\texttt{FinishResponse} (operation 15) and \texttt{ForceContinue}
(operation 16). There is no path to \texttt{spIdle} from
\texttt{spResponding} that bypasses the Stop hook. If a Stop hook
injects escalation guidance, the agent will see it.

\textbf{Gap:} The model does not track \emph{whether} the agent actually
escalated. This would require a new state variable
(\texttt{escalationPending : ZBOOL}) set when the agent identifies an
unresolved decision and cleared when it sends \texttt{/write}. Worth
adding if escalation compliance becomes a hard requirement.

\textbf{Proposed extension:}

\begin{verbatim}
  escalationPending : ZBOOL
  ...
  % In FinishResponse:
  escalationPending = ztrue ==> hookResult? = hdBlock
\end{verbatim}

This would make escalation a hard gate: the agent cannot stop until it
has sent the escalation message.

\subsubsection{UC8 --- Report Task Completion}

\textbf{Z coverage:} Partially modeled. The \texttt{PostToolUse} hook
on PR creation (existing \texttt{pr-announce.sh}) fires after
\texttt{CompleteToolSuccess}. The \texttt{Stop} hook at session end
maps to \texttt{FinishResponse}.

\textbf{What the model proves:} \texttt{CompleteToolSuccess} always
fires after successful tool execution, and the PostToolUse hook fires
as a side effect. There is no path where a PR is created but the
PostToolUse hook is skipped.

\textbf{Gap:} The model does not track whether a PR was created during
the session (\texttt{prCreated : ZBOOL}). Without this, the Stop hook
cannot conditionally remind about announcement. Worth adding for the
session close protocol validation.

\subsubsection{UC9 --- Redirect an Agent's Priorities}

\textbf{Z coverage:} Modeled via \texttt{ReceiveMessage}.

\textbf{What the model proves:} A message can arrive at any time during
an active session (\texttt{ReceiveMessage} fires in \texttt{spIdle},
\texttt{spProcessing}, or \texttt{spResponding}). The agent will
eventually see the redirect --- either via push notification (tool
description mutation, outside the model) or on its next \texttt{/read}.

\textbf{Gap:} The model does not capture the agent's response to
redirection (changing plan, switching tasks). This is behavioral, not
state-machine-level.

\subsubsection{UC10 --- Answer an Agent's Pending Question}

\textbf{Z coverage:} Same as UC9 --- modeled via \texttt{ReceiveMessage}
followed by \texttt{ReadMessages}.

\textbf{What the model proves:} The receive-then-read sequence is
always possible: \texttt{ReceiveMessage} increments the count, and
\texttt{ReadMessages} clears it. No deadlock between these operations.

\textbf{Gap:} None for V1.

\subsubsection{UC12 --- Coordinate File Ownership in a Shared Repo}

\textbf{Z coverage:} Partially modeled. The \texttt{SessionStart}
operation can inject coordination guidance (the \texttt{hasWall?} input
models context injection). \texttt{SetPlanManual} models the plan update
after negotiation.

\textbf{What the model proves:} The plan-before-work gate ensures that
both agents must declare their plans before editing files. If both agents
set plans, the plans are visible via \texttt{/who} (UC1), enabling
overlap detection. The model proves that \emph{no editing can begin
without a declared plan}, which is the foundation of file ownership
coordination.

\textbf{Gap:} The model does not enforce non-overlapping plans. Two
agents can both set plans and both edit files --- the model proves each
individually followed the protocol, not that the protocol prevented
conflicts. This matches the use case document's statement: ``biff is
coordination infrastructure, not a distributed lock.''

% ------------------------------------------------------------
\subsection{Group C: Broadcast}
% ------------------------------------------------------------

\subsubsection{UC6 --- Broadcast a Time-Sensitive Announcement}

\textbf{Z coverage:} Fully modeled.

\begin{itemize}
  \item \texttt{PostWall} --- sets \texttt{wallActive = ztrue},
    \texttt{wallText = \{text?\}}.
  \item \texttt{ClearWall} --- sets \texttt{wallActive = zfalse},
    \texttt{wallText = $\emptyset$}.
  \item \texttt{StartSession} --- loads wall state at session start
    via \texttt{hasWall?} and \texttt{wallContent?} inputs.
\end{itemize}

\textbf{What the model proves:}

\begin{enumerate}
  \item \textbf{Wall consistency}: \texttt{wallText} exists iff
    \texttt{wallActive = ztrue}. No ghost text, no active wall without
    content (invariant 3).
  \item \textbf{Wall visibility at startup}: \texttt{StartSession}
    propagates the wall into session state. A new session joining after
    a wall post will have \texttt{wallActive = ztrue} if a wall exists.
  \item \textbf{ClearWall requires active wall}: precondition
    \texttt{wallActive = ztrue} prevents clearing a non-existent wall.
\end{enumerate}

\textbf{Gap:} Wall TTL (expiry) is not modeled. The wall persists until
explicitly cleared. Adding TTL would require a time dimension
(monotonic clock), which significantly increases model complexity.

\subsubsection{UC15 --- Notify Team of CI/CD Status}

\textbf{Z coverage:} Modeled via \texttt{PostWall} (for \texttt{/wall})
and \texttt{ReceiveMessage} (for \texttt{/write} from GitHub to
committer).

\textbf{What the model proves:} The GitHub actor uses the same
operations as any other actor. The model does not distinguish actor
types (matching the use case document: ``The system does not distinguish
actor types at the protocol level''). All invariants hold regardless of
who triggers the operation.

\textbf{Gap:} GitHub's service token authentication is outside the
model (T2 configuration). The model assumes the actor is already
authenticated and has a valid session.

% ------------------------------------------------------------
\subsection{Group D: Real-Time Conversation}
% ------------------------------------------------------------

\subsubsection{UC14 --- Have a Synchronous Exchange}

\textbf{Z coverage:} Fully modeled.

\begin{itemize}
  \item \texttt{StartTalk} --- sets \texttt{talkActive = ztrue},
    guarded by \texttt{talkActive = zfalse} (no double-start).
  \item \texttt{EndTalk} --- sets \texttt{talkActive = zfalse},
    guarded by \texttt{talkActive = ztrue} (no double-end).
\end{itemize}

\textbf{What the model proves:}

\begin{enumerate}
  \item \textbf{Talk exclusivity}: exactly one talk session at a time.
    No path through the model allows \texttt{StartTalk} to fire twice
    without an intervening \texttt{EndTalk}.
  \item \textbf{Talk survives tool calls}: \texttt{talkActive} is
    preserved through the \texttt{BiffFrame} abbreviation. Tool calls,
    subagent spawns, and response cycles do not accidentally clear the
    talk state.
  \item \textbf{Clean shutdown}: \texttt{EndSession} clears
    \texttt{talkActive}, so no orphan talk sessions persist after logout.
\end{enumerate}

\textbf{Gap:} The model does not capture individual talk messages or
\texttt{/talk\_listen} blocking behavior. Talk messages are modeled the
same as \texttt{ReceiveMessage}. The blocking semantics of
\texttt{talk\_listen} (agent waits for message) would require a new tool
phase state, which adds complexity without clear workflow invariant value.

% ------------------------------------------------------------
\subsection{Group E: Focus Management}
% ------------------------------------------------------------

\subsubsection{UC13 --- Enter Do-Not-Disturb Mode}

\textbf{Z coverage:} Not modeled. The \texttt{mesgMode} variable was
proposed in the Layer 2 design but not included in the current spec to
keep the state space manageable.

\textbf{What the model proves:} Nothing specific to DND. However, the
\texttt{ReceiveMessage} operation fires regardless of DND state (matching
the use case: ``Messages sent to this actor still queue''). DND only
affects push notification suppression, which is a display concern outside
the state machine.

\textbf{Proposed extension (if needed):}

\begin{verbatim}
  mesgMode : ZBOOL  % ztrue = accepting, zfalse = DND
  ...
  % Invariant: walls bypass DND (already modeled ---
  % wallActive is independent of mesgMode)
\end{verbatim}

\textbf{Recommendation:} Do not add unless a workflow invariant depends
on DND state. Currently no use case requires ``the agent must be in
DND before doing X.''

% ------------------------------------------------------------
\subsection{Group F: Agent-to-Agent Coordination}
% ------------------------------------------------------------

\subsubsection{UC11 --- Avoid Duplicate Work Across Agents}

\textbf{Z coverage:} Partially modeled. The plan-before-work gate
ensures every agent declares work before editing. The
\texttt{StartSession} operation's wall injection models the SessionStart
hook's coordination guidance.

\textbf{What the model proves:} Every active editing session has
\texttt{planSet = ztrue}. This means \texttt{/who} output will always
show a plan for every working agent. The plan-before-work invariant is
the formal foundation for duplicate avoidance: if every agent must
declare intent before acting, the declarations are observable, and
conflicts are detectable.

\textbf{Gap:} The model does not verify that the agent actually checks
\texttt{/who} before claiming work. This is a behavioral directive
(injected via SessionStart hook), not a state machine constraint. To
formalize it, you would need a \texttt{whoChecked : ZBOOL} variable
set by a \texttt{/who} query and required as a precondition on
\texttt{ClaimBead}. This is feasible but may be over-constraining ---
the agent might legitimately claim a bead without checking \texttt{/who}
if it knows no other agents are active.

% ------------------------------------------------------------
\subsection{Coverage Summary}
% ------------------------------------------------------------

\begin{tabular}{llll}
\hline
\textbf{UC} & \textbf{Title} & \textbf{Coverage} & \textbf{Proven} \\
\hline
UC1 & Check Who's Working & Read-only & Plan state coherent \\
UC2 & Inspect Availability & Read-only & Talk state coherent \\
UC3 & Declare Plan & Full & Manual blocks auto; plan gates work \\
UC4 & Send Message & Side effect & No state corruption \\
UC5 & Check Messages & Full & Async receive + sync read correct \\
UC6 & Broadcast & Full & Wall consistency + startup visibility \\
UC7 & Escalate Decision & Hook point & Stop hook covers all paths \\
UC8 & Report Completion & Hook point & PostToolUse always fires \\
UC9 & Redirect Agent & Via receive & Messages arrive in any phase \\
UC10 & Answer Question & Via receive & Receive-read sequence sound \\
UC11 & Avoid Duplicate Work & Plan gate & All editors have plans \\
UC12 & Coordinate Ownership & Plan gate & Plan-before-work enforced \\
UC13 & DND Mode & Not modeled & No workflow invariant depends on it \\
UC14 & Synchronous Exchange & Full & Talk exclusivity + clean shutdown \\
UC15 & CI/CD Status & Via wall+msg & Actor-agnostic operations \\
\hline
\end{tabular}

\textbf{Legend:}
\begin{itemize}
  \item \textbf{Full} --- all state transitions modeled and verified.
  \item \textbf{Hook point} --- the Z spec proves the hook fires on all
    paths, but the hook's behavioral effect is not formally tracked.
  \item \textbf{Side effect} --- the action occurs but does not change
    tracked state.
  \item \textbf{Read-only} --- no state mutation; the model proves the
    queried state is consistent.
  \item \textbf{Not modeled} --- deliberately excluded; no workflow
    invariant requires it.
\end{itemize}

\end{document}
